\section{General concentration and counting estimates}
\label{sec:general-estimates}\label{sec:layer1}

Let $(a_k)_{k\geq0}$ be a sequence of random complex coefficients on one
probability space, and put
\[
 f_N(z)=\sum_{k=0}^N a_kz^k.
\]
The same coefficient sequence is used for all degrees, and in particular
$\deg f_N\leq N$ when $f_N$ is nonzero.
Let $\sigma_N(r)>0$ be deterministic, and write
\[
 J_N(r)=\int_{\mathbb T}\log|f_N(re^{i\theta})|\dd\mu(\theta),
\]
where $\mu$ is the normalized angular measure. Deterministic counting
statements concern nonzero polynomials, and in probabilistic
applications we include the case of an identically zero polynomial in an
explicit exceptional event.

\subsection{Angular occupation and logarithmic integrals}

\begin{namedhypothesis}[One- and two-point occupation bounds]\phantomsection\label{hyp:occupation}
For fixed $N,r$, we set
\[
 X(\theta)=\frac{f_N(re^{i\theta})}{\sigma_N(r)}.
\]
Let $G$ have density $\pi^{-1}e^{-|z|^2}$ on $\C$. We set
\[
 F_G(t)=\Prob(|G|\leq t)=1-e^{-t^2}.
\]
There are measurable exceptional sets
$E_1\subset\mathbb T$, $E_2\subset\mathbb T^2$ with
$\mu(E_1)\leq\eta_1$, $(\mu\otimes\mu)(E_2)\leq\eta_2$, such that, for
all $s,t\geq0$,
\[
 \left|\Prob(|X(\theta)|\leq t)-F_G(t)\right|\leq\delta_1
       \quad(\theta\notin E_1),
\]
\[
 \left|\Prob(|X(\theta)|\leq s,|X(\phi)|\leq t)-F_G(s)F_G(t)\right|
       \leq\delta_2\quad((\theta,\phi)\notin E_2).
\]
\end{namedhypothesis}
In the two-point bound we compare with a pair of independent standard
complex Gaussians, and if the normal approximation yields a correlated
Gaussian pair, then $\delta_2$ includes the error of the comparison with
this independent pair.

\begin{namedhypothesis}[Logarithmic moments and energy]\phantomsection\label{hyp:logenergy}
For fixed constants $C>0$, $\beta>0$, and for all real $p>1$,
\begin{equation}
 \E\int_{\mathbb T}|\log|X(\theta)||^p\dd\mu(\theta)
       \leq(Cp)^{\beta p}. \label{eq:log-moment-bound}
\end{equation}
We will also use the following eventwise bound on an event
$E_H$ with $\Prob(E_H^c)\leq p_H$:
\begin{equation}
 \E\left[\boldsymbol1_{E_H}
       \int_{\mathbb T}|\log|X(\theta)||^p\dd\mu(\theta)\right]
       \leq(Cp)^{\beta p}\qquad(p>1). \label{eq:event-log-moment-bound}
\end{equation}
Here the expectation is the integral over $E_H$. We evaluate the
integrand only on that event. This also defines the eventwise bound when the angular
integral is infinite on $E_H^c$.
The exact angular energy identity is
\begin{equation}
 \int_{\mathbb T}|X(\theta)|^2\dd\mu(\theta)=1
       \quad\hbox{almost surely}. \label{eq:exact-angular-energy}
\end{equation}
The upper-energy probability bound is
\begin{equation}
 \Prob\left(\int_{\mathbb T}|X(\theta)|^2\dd\mu(\theta)>B_E\right)
       \leq p_E \qquad(B_E>0). \label{eq:angular-energy-tail}
\end{equation}
\end{namedhypothesis}
The full bound \eqref{eq:log-moment-bound} is the case
$E_H=\Omega$ and $p_H=0$ of \eqref{eq:event-log-moment-bound}.
We use the tail form \eqref{eq:angular-energy-tail} when the coefficient
moduli are random, and under \eqref{eq:exact-angular-energy} we can take
$B_E=1$, $p_E=0$.
For the coefficient distributions of Sections~\ref{sec:concentration}
and~\ref{sec:layer2-models}, these logarithmic bounds hold for every
$p\ge1$ (see \eqref{eq:layer2-nns}, \eqref{eq:layer2-gaussianlog} and
Proposition~\ref{prop:coefficient-transfer}), but in the proof of
Theorem~\ref{thm:T1} we use them only for $p=2q\ge2$.

\begin{theorem}[Logarithmic integral concentration]\label{thm:T1}
Assume the \hyperref[hyp:occupation]{one- and two-point occupation bounds},
the eventwise logarithmic bound \eqref{eq:event-log-moment-bound},
and the energy bound \eqref{eq:angular-energy-tail}. We set
\[
 d_1=\delta_1+\eta_1,\qquad d_2=\delta_2+\eta_2,
 \qquad v=d_2+2d_1.
\]
For all $T>0$, $u>0$, $b>0$, $q\geq1$, and $\lambda>1$, we set
\begin{align}
 \eta_{\log}={}&u+2Td_1+\lambda(2Cq)^\beta
        \sqrt{e^{-2T}+d_1+b}
        +(B_E/2+1)e^{-2T}, \label{eq:T1-error}\\
 p_{\rm fail}={}&v/b^2+\lambda^{-2q}+4T^2v/u^2+p_E+p_H. \label{eq:T1-failure}
\end{align}
Then
\[
 \Prob\left(\left|J_N(r)-\log\sigma_N(r)
                   -\E\log|G|\right|>\eta_{\log}\right)\leq p_{\rm fail},
 \qquad \E\log|G|=-\gamma/2.
\]
The same conclusion holds if, in place of the occupation bounds,
$|\E A_t-F_G(t)|\leq d_1$ and $\operatorname{Var}(A_t)\leq v$ for all
$t\geq0$ and some $d_1,v\geq0$, where
$A_t=\mu\{\theta:|X(\theta)|\leq t\}$. If moreover
$\int_{\mathbb T}|X|^2\dd\mu\leq B_E$ on $E_H$, then $p_E+p_H$ in
\eqref{eq:T1-failure} may be replaced by $\Prob(E_H^c)$.
\end{theorem}

\begin{proof}
The idea is to clip the logarithm at $\pm T$, to estimate the mean and
the variance of the clipped integral by the occupation bounds, and to
control the lower tail that the clipping removes by the logarithmic
moments and the occupation bound at $e^{-T}$, and the upper tail by the
energy bound. Table~\ref{tab:T1-events} lists the five exceptional events,
one for each term of \eqref{eq:T1-failure}.

We set
$A_t=\int_{\mathbb T}\boldsymbol1_{\{|X(\theta)|\leq t\}}\dd\mu(\theta)$.
By Tonelli's theorem, we have
\begin{align*}
 \E A_t
 &=\int_{\mathbb T}\Prob(|X(\theta)|\leq t)\dd\mu(\theta),\\
 \E A_t^2
 &=\E\int_{\mathbb T^2}
      \boldsymbol1_{\{|X(\theta)|\leq t\}}
      \boldsymbol1_{\{|X(\phi)|\leq t\}}\dd\mu(\theta)\dd\mu(\phi)\\
 &=\int_{\mathbb T^2}
      \Prob(|X(\theta)|\leq t,|X(\phi)|\leq t)
         \dd\mu(\theta)\dd\mu(\phi).
\end{align*}
By these identities and the
\hyperref[hyp:occupation]{one- and two-point occupation bounds}, with $s=t$
in the second, and since the probabilities and $F_G(t)$ lie in $[0,1]$, we
have, for every $t\geq0$,
\begin{align*}
 |\E A_t-F_G(t)|
 &=\left|\int_{\mathbb T}\bigl(\Prob(|X(\theta)|\leq t)-F_G(t)\bigr)
       \dd\mu(\theta)\right|\\
 &\leq\int_{\mathbb T\setminus E_1}\delta_1\dd\mu
       +\int_{E_1}1\dd\mu
 \leq\delta_1+\mu(E_1)
 \leq d_1,\\
 |\E A_t^2-F_G(t)^2|
 &=\left|\int_{\mathbb T^2}\bigl(\Prob(|X(\theta)|\leq t,|X(\phi)|\leq t)
       -F_G(t)^2\bigr)\dd(\mu\otimes\mu)\right|\\
 &\leq\int_{\mathbb T^2\setminus E_2}\delta_2\dd(\mu\otimes\mu)
       +\int_{E_2}1\dd(\mu\otimes\mu)
 \leq\delta_2+(\mu\otimes\mu)(E_2)
 \leq d_2.
\end{align*}
Since both $\E A_t$ and $F_G(t)$ belong to $[0,1]$, we have
\begin{gather*}
 |(\E A_t)^2-F_G(t)^2|
 =|\E A_t-F_G(t)|\,|\E A_t+F_G(t)|
 \leq2d_1,\\
 \operatorname{Var}(A_t)
 =\E A_t^2-(\E A_t)^2
 \leq|\E A_t^2-F_G(t)^2|+|(\E A_t)^2-F_G(t)^2|
 \leq d_2+2d_1
 =v.
\end{gather*}
From here on, the occupation bounds enter only through the bounds
$|\E A_t-F_G(t)|\leq d_1$ and $\operatorname{Var}(A_t)\leq v$, which
proves the first variant in the statement.

Since $F_G(e^{-T})=1-\exp(-e^{-2T})\leq e^{-2T}$, taking $t=e^{-T}$ in
these bounds and applying Chebyshev's inequality, we obtain
\begin{align*}
 \E A_{e^{-T}}
 &\leq F_G(e^{-T})+d_1
 \leq e^{-2T}+d_1,\\
 \Prob(A_{e^{-T}}>e^{-2T}+d_1+b)
 &\leq\Prob(A_{e^{-T}}-\E A_{e^{-T}}>b)
 \leq\frac{\operatorname{Var}(A_{e^{-T}})}{b^2}
 \leq\frac v{b^2}.
\end{align*}
We call the event $\{A_{e^{-T}}>e^{-2T}+d_1+b\}$ the Chebyshev event for
$A_{e^{-T}}$.

On $E_H$, we set
$U=\int_{\mathbb T}|\log|X(\theta)||^2\dd\mu(\theta)$.
By Jensen's inequality and \eqref{eq:event-log-moment-bound} with $p=2q$,
we have
\begin{align*}
 U^q
 &=\left(\int_{\mathbb T}|\log|X(\theta)||^2\dd\mu(\theta)\right)^q
 \leq\int_{\mathbb T}|\log|X(\theta)||^{2q}\dd\mu(\theta),\\
 \E[\boldsymbol1_{E_H}U^q]
 &\leq\E\left[\boldsymbol1_{E_H}
       \int_{\mathbb T}|\log|X(\theta)||^{2q}\dd\mu(\theta)\right]
 \leq(2Cq)^{2\beta q}.
\end{align*}
Applying Markov's inequality to $U^q$ on $E_H$, we get
\[
 \Prob\{\omega\in E_H:U>\lambda^2(2Cq)^{2\beta}\}
 \leq\frac{\E[\boldsymbol1_{E_H}U^q]}
              {\lambda^{2q}(2Cq)^{2\beta q}}
 \leq\frac{(2Cq)^{2\beta q}}{\lambda^{2q}(2Cq)^{2\beta q}}
 =\lambda^{-2q}.
\]
We call the event on the left the Markov event for $U$.

We now truncate the logarithm and set
$\ell_T(z)=\max(-T,\min(T,\log|z|))$ and
$L_T=\int_{\mathbb T}\ell_T(X)\dd\mu$.
For every $z\in\C$, including $z=0$, we have the layer-cake identity
\[
 \ell_T(z)=T-\int_{-T}^T\boldsymbol1_{\{|z|\leq e^\tau\}}\dd\tau,
\]
since the set of $\tau\in[-T,T]$ with $|z|\leq e^\tau$ is
$[\max(-T,\log|z|),T]$ if $\log|z|\leq T$ (with $\log0=-\infty$), and is
empty otherwise.
By the layer-cake identity at $z=X(\theta)$ and at $z=G$, Tonelli's
theorem and $\Prob(|G|\leq e^\tau)=F_G(e^\tau)$, we have
\begin{align*}
 L_T
 &=\int_{\mathbb T}\left(T-\int_{-T}^T
         \boldsymbol1_{\{|X(\theta)|\leq e^\tau\}}\dd\tau\right)\dd\mu(\theta)
 =T-\int_{-T}^T A_{e^\tau}\dd\tau,\\
 \E L_T-\E\ell_T(G)
 &=\int_{-T}^T\bigl(F_G(e^\tau)-\E A_{e^\tau}\bigr)\dd\tau,\\
 L_T-\E L_T
 &=-\int_{-T}^T(A_{e^\tau}-\E A_{e^\tau})\dd\tau.
\end{align*}
By these identities, the bounds $|\E A_t-F_G(t)|\leq d_1$ and
$\operatorname{Var}(A_t)\leq v$, and the Cauchy--Schwarz inequality, we
have
\begin{align*}
 |\E L_T-\E\ell_T(G)|
 &\leq\int_{-T}^T|\E A_{e^\tau}-F_G(e^\tau)|\dd\tau
 \leq2Td_1,\\
 |\operatorname{Cov}(A_{e^\tau},A_{e^{\tau'}})|
 &\leq\sqrt{\operatorname{Var}(A_{e^\tau})\operatorname{Var}(A_{e^{\tau'}})}
 \leq v,
\end{align*}
and hence, by Fubini's theorem, which applies since
$|A_{e^\tau}-\E A_{e^\tau}|\leq1$,
\begin{align*}
 &\operatorname{Var}(L_T)
 =\E\left(\int_{-T}^T(A_{e^\tau}-\E A_{e^\tau})\dd\tau\right)^2\\
 &\quad=\int_{-T}^T\int_{-T}^T
       \E[(A_{e^\tau}-\E A_{e^\tau})(A_{e^{\tau'}}-\E A_{e^{\tau'}})]
       \dd\tau\dd\tau'\\
 &\quad=\int_{-T}^T\int_{-T}^T
       \operatorname{Cov}(A_{e^\tau},A_{e^{\tau'}})\dd\tau\dd\tau'
 \leq\int_{-T}^T\int_{-T}^T
       |\operatorname{Cov}(A_{e^\tau},A_{e^{\tau'}})|\dd\tau\dd\tau'
 \leq4T^2v.
\end{align*}
Applying Chebyshev's inequality again, we get
\[
 \Prob(|L_T-\E L_T|>u)
 \leq\frac{\operatorname{Var}(L_T)}{u^2}
 \leq\frac{4T^2v}{u^2}.
\]
We call the event on the left the Chebyshev event for $L_T$.

On $E_H$, outside the Chebyshev event for $A_{e^{-T}}$ and the Markov
event for $U$, by the Cauchy--Schwarz inequality we bound the negative
truncation error by
\begin{align*}
 \int_{\{|X|<e^{-T}\}}(-\log|X|-T)\dd\mu
 &\leq\int_{\{|X|<e^{-T}\}}|\log|X||\dd\mu\\
 &\leq\left(\int_{\mathbb T}|\log|X||^2\dd\mu\right)^{1/2}
       \left(\int_{\mathbb T}
         \boldsymbol1_{\{|X|<e^{-T}\}}\dd\mu\right)^{1/2}\\
 &\leq\sqrt{UA_{e^{-T}}}
 \leq\lambda(2Cq)^\beta\sqrt{e^{-2T}+d_1+b}.
\end{align*}
On the energy event $\{\int_{\mathbb T}|X|^2\dd\mu\leq B_E\}$, by
Tonelli's theorem applied to
$(y-T)_+=\int_T^\infty\boldsymbol1_{\{y>\tau\}}\dd\tau$ at $y=\log|X|$
and by Markov's inequality for $\mu$, we bound the positive truncation
error by
\begin{align*}
 \int_{\{|X|>e^T\}}(\log|X|-T)\dd\mu
 &=\int_T^\infty\mu\{|X|>e^\tau\}\dd\tau
 \leq\int_T^\infty e^{-2\tau}\Bigl(\int_{\mathbb T}|X|^2\dd\mu\Bigr)\dd\tau\\
 &\leq\int_T^\infty B_E e^{-2\tau}\dd\tau
 =\frac{B_E}2e^{-2T}.
\end{align*}
For the Gaussian tails,
$\int_T^\infty\Prob(|G|<e^{-\tau})\dd\tau$ and
$\int_T^\infty\Prob(|G|>e^\tau)\dd\tau$ are at most
$\frac12e^{-2T}$ by \eqref{eq:gaussian-log-tails}.
On $E_H$, outside the Markov event for $U$, the integral
$\int_{\mathbb T}\log|X|\dd\mu$ is finite, since
$\int_{\mathbb T}|\log|X||^2\dd\mu=U\leq\lambda^2(2Cq)^{2\beta}$. Since
\[
 |\ell_T(z)-\log|z||=(-\log|z|-T)_++(\log|z|-T)_+,
\]
on $E_H$ and on the energy event, outside the Chebyshev event for
$A_{e^{-T}}$ and the Markov event for $U$, it follows from the layer-cake
identity, the two truncation bounds and the two Gaussian tail bounds that
\begin{align*}
 \left|\int_{\mathbb T}\log|X|\dd\mu-L_T\right|
 &\leq\int_{\mathbb T}|\log|X|-\ell_T(X)|\dd\mu\\
 &=\int_{\{|X|<e^{-T}\}}(-\log|X|-T)\dd\mu
      +\int_{\{|X|>e^T\}}(\log|X|-T)\dd\mu\\
 &\leq\lambda(2Cq)^\beta\sqrt{e^{-2T}+d_1+b}+\frac{B_E}2e^{-2T},\\
 |\E\ell_T(G)-\E\log|G||
 &\leq\E|\ell_T(G)-\log|G||
 =\E(-\log|G|-T)_++\E(\log|G|-T)_+\\
 &=\int_T^\infty\Prob(|G|<e^{-\tau})\dd\tau
       +\int_T^\infty\Prob(|G|>e^\tau)\dd\tau
 \leq e^{-2T}.
\end{align*}

\begin{table}[!htbp]
\centering
\begin{tabular}{@{}lll@{}}
\toprule
Exceptional event & Defined by & Probability at most\\
\midrule
Chebyshev event for $A_{e^{-T}}$ & $A_{e^{-T}}>e^{-2T}+d_1+b$ & $v/b^2$\\
Markov event for $U$ & $\omega\in E_H$, $U>\lambda^2(2Cq)^{2\beta}$
 & $\lambda^{-2q}$\\
Chebyshev event for $L_T$ & $|L_T-\E L_T|>u$ & $4T^2v/u^2$\\
Complement of the energy event & $\int_{\mathbb T}|X|^2\dd\mu>B_E$
 & $p_E$\\
Complement of $E_H$ & & $p_H$\\
\bottomrule
\end{tabular}
\caption{The exceptional events in the proof of Theorem~\ref{thm:T1}.
The sum of the bounds on their probabilities is
\eqref{eq:T1-failure}.}
\label{tab:T1-events}
\end{table}
Finally, $J_N(r)-\log\sigma_N(r)=\int\log|X|\dd\mu$, and outside the five
exceptional events of Table~\ref{tab:T1-events}, by the previous display
and the bounds on $|L_T-\E L_T|$ and $|\E L_T-\E\ell_T(G)|$, we have
\begin{align*}
 &\left|\int\log|X|\dd\mu-\E\log|G|\right|
 \leq\left|\int\log|X|\dd\mu-L_T\right|
       +|L_T-\E L_T|\\
 &\qquad+|\E L_T-\E\ell_T(G)|
       +|\E\ell_T(G)-\E\log|G||\\
 &\qquad\leq\lambda(2Cq)^\beta\sqrt{e^{-2T}+d_1+b}+\frac{B_E}2e^{-2T}
       +u+2Td_1+e^{-2T}
 =\eta_{\log}.
\end{align*}
Under \eqref{eq:exact-angular-energy} and \eqref{eq:log-moment-bound} the
last two events are absent. If $\int_{\mathbb T}|X|^2\dd\mu\leq B_E$ on $E_H$,
then the complement of the energy event is contained in $E_H^c$, so that
the last two events have probability at most $\Prob(E_H^c)$ together,
which proves the second variant in the statement. By \eqref{eq:event-log-moment-bound},
the logarithmic integral is finite almost surely on $E_H$, even for
polynomials with roots on the integration circle.
Since $|G|^2$ has the exponential distribution with mean one,
\eqref{eq:euler-centering} gives $\E\log|G|=-\gamma/2$.
\end{proof}

\begin{remark}[Polynomial choices of parameters]\label{rem:T1-powers}
If $d_1\leq D N^{-a}$ and $v\leq V N^{-a}$, choose
$T=t\log N$, $b=N^{-2t}$, $u=N^{-w}$, and $q=c\log N\geq1$.
Under \eqref{eq:log-moment-bound} and \eqref{eq:exact-angular-energy},
we may take $p_H=0$, $B_E=1$ and $p_E=0$. With $\lambda=e$, so that
$\lambda(2Cq)^\beta=e(2Cc\log N)^\beta$, the terms of
\eqref{eq:T1-error} and \eqref{eq:T1-failure}, other than $u=N^{-w}$ and
$p_E=p_H=0$, satisfy
\begin{align*}
 2Td_1&\leq2t(\log N)DN^{-a},\qquad
 (B_E/2+1)e^{-2T}=\tfrac32N^{-2t},\\
 e^{-2T}+d_1+b&\leq N^{-2t}+DN^{-a}+N^{-2t}=2N^{-2t}+DN^{-a},\\
 v/b^2&\leq VN^{-a}/N^{-4t}=VN^{4t-a},\qquad
 \lambda^{-2q}=e^{-2c\log N}=N^{-2c},\\
 4T^2v/u^2&\leq4t^2(\log N)^2VN^{-a}/N^{-2w}=4t^2V(\log N)^2N^{2w-a}.
\end{align*}
By substituting these bounds, we get
\[
 \eta_{\log}\leq N^{-w}+2tD(\log N)N^{-a}
  +e(2Cc\log N)^\beta\sqrt{2N^{-2t}+DN^{-a}}
  +\tfrac32N^{-2t},
\]
\[
 p_{\rm fail}\leq VN^{4t-a}+N^{-2c}
       +4t^2V(\log N)^2N^{2w-a}.
\]
The failure terms tend to zero when $4t<a$ and $2w<a$.
For $N=j^s$, the three failure terms are
\[
 Vj^{-s(a-4t)},\qquad j^{-2cs},\qquad
 4t^2Vs^2(\log j)^2j^{-s(a-2w)},
\]
and they are summable if
\[
 s(a-4t)>1,\qquad 2cs>1,\qquad s(a-2w)>1.
\]
In particular, take $a=1/2$, $t=w=1/32$, $c=6$, $\beta=6$ and $s=8$.
Then $s(a-4t)=8(\frac12-\frac18)=3$, $2cs=2\cdot6\cdot8=96$ and
$s(a-2w)=8(\frac12-\frac1{16})=\frac72$ all exceed one, and for $N\geq e$,
\begin{align*}
 \eta_{\log}
 &\leq N^{-1/32}+\frac D{16}(\log N)N^{-1/2}
  +e(12C\log N)^6\sqrt{2N^{-1/16}+DN^{-1/2}}
  +\tfrac32N^{-1/16}\\
 &\leq\Bigl(\frac52+\frac D{16}+e(12C)^6\sqrt{2+D}\Bigr)
       N^{-1/32}(\log N)^6,
\end{align*}
by the bound on $\eta_{\log}$ above with $2tD=D/16$, $2Cc=12C$,
$2t=1/16$ and $w=1/32$, and by $\log N\geq1$ and
$N^{-1/2}\leq N^{-1/16}\leq N^{-1/32}$.
So we get an error bounded by the above constant times
$N^{-1/32}(\log N)^6$ and a summable failure on $j^8$. For large enough
degree the error is at most $N^{-1/32}(\log N)^7$, with a deterministic
threshold depending on the fixed constant: it suffices that $\log N$ is at
least the above constant. We use the moments of the angular
logarithmic norm in \eqref{eq:log-moment-bound}.
The proof of Lemma~\ref{lem:clipping} takes instead
$u=N^{-1/32}(\log N)^7$, $q=\lceil\log N\rceil$ and
$\lambda=(2e\log N/q)^6\log N\ge e^6\log N$,
for which $\lambda(2Cq)^6=(4eC)^6(\log N)^7$ and
$\lambda^{-2q}\le N^{-12}$. The choice $q=6\log N$, $\lambda=e$ above
gives the same failure bound $N^{-2c}=N^{-12}$ with
$\lambda(2Cq)^6=e(12C\log N)^6$. The extra factor $\log N$ in $\lambda$
is the difference between the exponents $7$ and $6$.
\end{remark}

\subsection{A local radial secant estimate}

\begin{lemma}[Local approximation of radial counts]\label{lem:local-radial-secant}
Let $f_N$ be a nonzero polynomial of degree at most $N$, and let
$\sigma_N(r)>0$ be deterministic. Suppose that the function $F_N(x)$,
equal to $\log\sigma_N(1+x/N)$ minus a centering independent of $x$,
satisfies $|F_N(x)-F(x)|\leq\varepsilon$ on $[-K,K]$.
Suppose also $|F'|\leq B_F$ and
$|F'(x)-F'(y)|\leq M_F|x-y|$ there.
Suppose the three logarithmic deviations
\[
 \left|J_N(1+y/N)-\log\sigma_N(1+y/N)-\E\log|G|\right|,
 \qquad y=x-h,x,x+h,
\]
are at most $\eta_{\log}$ outside their respective exceptional events.
If $h>0$ and $|x|+h\leq K<N$, then outside the union of those three events,
\begin{equation}
 \left|\frac{\nu_N(1+x/N)}N-F'(x)\right|
 \leq\frac{M_Fh}{2}+\frac{B_FK}{N}
       +\frac{2(\eta_{\log}+\varepsilon)(1+K/N)}h. \label{eq:T2-moving}
\end{equation}
In particular an
error $\eta_{\log}+\varepsilon\leq N^{-1/32}(\log N)^7$, with $h=N^{-\kappa}$,
$0<\kappa\leq1/64$, gives convergence to $F'(0)$ uniformly for
deterministic $|x|\leq h$.
\end{lemma}

\begin{proof}
By Jensen's formula, $\nu_N(r)/N$ lies between two secant slopes of
$J_N$ in the variable $N\log\rho$, one on each side of $r$, and the idea
is to replace $J_N$ by $F$ at the three probe points, up to a constant
which cancels, and then to compare the secant slopes of $F$ with $F'(x)$.
This leads to three errors, which come from the curvature of $F$, the
change of the denominators, and the deviations of $J_N$ from $F$ at the
probe points.

Let $r=1+x/N$ and set
\[
 d_-=N\log\frac{r}{r-h/N},\qquad
 d_+=N\log\frac{r+h/N}{r}.
\]
Since $|x|+h\leq K<N$, we have $r-h/N=1+(x-h)/N\geq1-K/N>0$.
By the integral form \eqref{eq:jensen-integral} of Jensen's formula,
in which $\nu_N$ counts the roots in the closed disk, and since $\nu_N$ is
nondecreasing, we have
\begin{align*}
 J_N(r)-J_N(r-h/N)
 &=\int_{r-h/N}^r\frac{\nu_N(\rho)}\rho\dd\rho
 \leq\nu_N(r)\log\frac r{r-h/N}
 =\frac{\nu_N(r)d_-}N,\\
 J_N(r+h/N)-J_N(r)
 &=\int_r^{r+h/N}\frac{\nu_N(\rho)}\rho\dd\rho
 \geq\nu_N(r)\log\frac{r+h/N}r
 =\frac{\nu_N(r)d_+}N.
\end{align*}
Dividing by $d_\pm>0$, we obtain the Jensen bounds
\[
 \frac{J_N(r)-J_N(r-h/N)}{d_-}
 \leq\frac{\nu_N(r)}N
 \leq\frac{J_N(r+h/N)-J_N(r)}{d_+}.
\]
Since $N\log(1+u/N)$ has derivative $1/(1+u/N)$ in $u$, we have
\begin{align*}
 d_-&=N\log(1+x/N)-N\log\bigl(1+(x-h)/N\bigr)
 =\int_{x-h}^x\frac{\mathrm du}{1+u/N},\\
 d_+&=N\log\bigl(1+(x+h)/N\bigr)-N\log(1+x/N)
 =\int_x^{x+h}\frac{\mathrm du}{1+u/N}.
\end{align*}
Since $|x|+h\leq K<N$, every $u\in[x-h,x+h]$ satisfies
$0<1-K/N\leq1+u/N\leq1+K/N$, and taking reciprocals and integrating over
an interval of length $h$, we obtain
\begin{gather*}
 \frac1{1+K/N}\leq\frac1{1+u/N}\leq\frac1{1-K/N},\qquad
 \frac h{1+K/N}\leq d_\pm\leq\frac h{1-K/N},\\
 1-\frac KN\leq\frac h{d_\pm}\leq1+\frac KN,\qquad
 \left|\frac h{d_\pm}-1\right|\leq\frac KN.
\end{gather*}

Since $F'$ is Lipschitz on $[-K,K]$, which contains $x\pm t$ for
$0\leq t\leq h$, for the ordinary secants we have
\begin{align*}
 \left|\frac{F(x+h)-F(x)}h-F'(x)\right|
 &=\left|\frac1h\int_0^h\bigl(F'(x+t)-F'(x)\bigr)\dd t\right|\\
 &\leq\frac1h\int_0^h M_Ft\dd t
 =\frac{M_Fh}2,\\
 \left|\frac{F(x)-F(x-h)}h-F'(x)\right|
 &=\left|\frac1h\int_0^h\bigl(F'(x-t)-F'(x)\bigr)\dd t\right|\\
 &\leq\frac1h\int_0^h M_Ft\dd t
 =\frac{M_Fh}2.
\end{align*}
Each ordinary secant has absolute value at most $B_F$, since it is the
mean of $F'$ over an interval in $[-K,K]$.
Replacing its denominator $h$ by $d_\pm$ multiplies it by $h/d_\pm$, and
\begin{align*}
 \left|\frac{F(x+h)-F(x)}{d_+}-\frac{F(x+h)-F(x)}h\right|
 &=\left|\frac{F(x+h)-F(x)}h\right|\,\left|\frac h{d_+}-1\right|\\
 &\leq B_F\left|\frac h{d_+}-1\right|
 \leq\frac{B_FK}N,\\
 \left|\frac{F(x)-F(x-h)}{d_-}-\frac{F(x)-F(x-h)}h\right|
 &=\left|\frac{F(x)-F(x-h)}h\right|\,\left|\frac h{d_-}-1\right|\\
 &\leq B_F\left|\frac h{d_-}-1\right|
 \leq\frac{B_FK}N.
\end{align*}
Let $c$ be the centering, so that $F_N(x)=\log\sigma_N(1+x/N)-c$.
For $y=x-h,x,x+h$, which lie in $[-K,K]$, we have
\begin{align*}
 &J_N(1+y/N)-\E\log|G|-c-F(y)\\
 &\quad=\bigl\{J_N(1+y/N)-\log\sigma_N(1+y/N)-\E\log|G|\bigr\}
        +\bigl\{F_N(y)-F(y)\bigr\}.
\end{align*}
The constant $\E\log|G|+c$ does not depend on $y$, so that it cancels from
each secant, and outside the exceptional event at $y$ the first brace is
at most $\eta_{\log}$ in absolute value and the second is at most
$\varepsilon$. Since
$r\pm h/N=1+(x\pm h)/N$, by subtracting the identity at two neighboring
probe points we obtain
\begin{align*}
 \bigl|J_N(r+h/N)-J_N(r)-\bigl(F(x+h)-F(x)\bigr)\bigr|
 &\leq2(\eta_{\log}+\varepsilon),\\
 \bigl|J_N(r)-J_N(r-h/N)-\bigl(F(x)-F(x-h)\bigr)\bigr|
 &\leq2(\eta_{\log}+\varepsilon).
\end{align*}
By the lower bound $d_\pm\geq h/(1+K/N)$, the secant error is therefore at
most $2(\eta_{\log}+\varepsilon)/d_\pm\leq2(\eta_{\log}+\varepsilon)(1+K/N)/h$.
Adding the three errors to the Jensen bounds, we obtain
\begin{align*}
 \frac{\nu_N(r)}N
 &\leq\frac{J_N(r+h/N)-J_N(r)}{d_+}
 \leq\frac{F(x+h)-F(x)}{d_+}+\frac{2(\eta_{\log}+\varepsilon)(1+K/N)}h\\
 &\leq\frac{F(x+h)-F(x)}h+\frac{B_FK}N
      +\frac{2(\eta_{\log}+\varepsilon)(1+K/N)}h\\
 &\leq F'(x)+\frac{M_Fh}2+\frac{B_FK}N
      +\frac{2(\eta_{\log}+\varepsilon)(1+K/N)}h,\\
 \frac{\nu_N(r)}N
 &\geq\frac{J_N(r)-J_N(r-h/N)}{d_-}
 \geq\frac{F(x)-F(x-h)}{d_-}-\frac{2(\eta_{\log}+\varepsilon)(1+K/N)}h\\
 &\geq\frac{F(x)-F(x-h)}h-\frac{B_FK}N
      -\frac{2(\eta_{\log}+\varepsilon)(1+K/N)}h\\
 &\geq F'(x)-\frac{M_Fh}2-\frac{B_FK}N
      -\frac{2(\eta_{\log}+\varepsilon)(1+K/N)}h,
\end{align*}
which together give \eqref{eq:T2-moving}.
The Jensen bounds include roots on the counting circle.
Since the three exceptional events may be arbitrarily dependent, we add
their three probabilities.
To prove the last assertion, let $|x|\leq h=N^{-\kappa}$ and let $N$ be so
large that $2h\leq K$, so that $|x|+h\leq K$. Since
$|F'(x)-F'(0)|\leq M_F|x|\leq M_Fh$, it follows from
\eqref{eq:T2-moving} and $\eta_{\log}+\varepsilon\leq N^{-1/32}(\log N)^7$
that
\begin{align*}
 \left|\frac{\nu_N(1+x/N)}N-F'(0)\right|
 &\leq\left|\frac{\nu_N(1+x/N)}N-F'(x)\right|+|F'(x)-F'(0)|\\
 &\leq\frac{M_Fh}{2}+\frac{B_FK}{N}
       +\frac{2(\eta_{\log}+\varepsilon)(1+K/N)}h+M_Fh\\
 &=\frac{3M_F}{2}N^{-\kappa}+\frac{B_FK}{N}
       +2\Bigl(1+\frac KN\Bigr)(\eta_{\log}+\varepsilon)N^{\kappa}\\
 &\leq\frac{3M_F}{2}N^{-\kappa}+\frac{B_FK}{N}
       +2\Bigl(1+\frac KN\Bigr)N^{\kappa-1/32}(\log N)^7,
\end{align*}
which tends to zero since $\kappa\leq1/64<1/32$.
\end{proof}

\subsection{A second-moment bound for small-derivative roots}

\begin{namedhypothesis}[Mesh and local multiplicity]\phantomsection\label{hyp:mesh}
Let $\mathcal A$ be an annulus contained in $\{|z|\leq R\}$ and let
$z_1,\ldots,z_m$ be retained mesh points.  Apart from at most $b_0$
roots in explicitly excluded sectors, every root of $f_N$ in $\mathcal A$
is within distance $\tau$ of a retained point. The segment between them
lies in a specified domain $D$, and $|z_i|\leq R$.
Except on events of probabilities at most $p_2$ and $p_0$, respectively,
\[
 \sup_D|f_N''|\leq M_2,
 \qquad\hbox{each mesh cell contains at most $q_0$ roots.}
\]
The probability $p_0$ also
includes failure of the sector bound.
\end{namedhypothesis}
For the cell bound we can use a disk count, provided that each cell lies
in one of the disks of a local count such as
Lemma~\ref{lem:local-counts-general}.

\begin{namedhypothesis}[Jet approximation]\phantomsection\label{hyp:jets}
We write the pair
\[
 W(z_i)=(W_0(z_i),W_1(z_i))=(f_N(z_i)/\sqrt N,z_if_N'(z_i)/N^{3/2})
\]
as a real four-vector.
For every retained $i$ there is a mean-zero Gaussian four-vector
$H_i=((H_i)_0,(H_i)_1)$, written as a pair like $W(z_i)$,
with covariance least eigenvalue at least $\lambda>0$.  For the
product-of-disks event
$E_i=\{|W_0(z_i)|\leq u,|W_1(z_i)|\leq v_0\}$, put
$g_i=\Prob(|(H_i)_0|\leq u,|(H_i)_1|\leq v_0)$ and assume
\[
 |\Prob(E_i)-g_i|\leq\delta.
\]
All but at most $b_*m^2$ \emph{ordered} pairs $(i,j)$ additionally obey
\[
 |\Prob(E_i\cap E_j)-g_i g_j|\leq\delta_*+\chi.
\]
\end{namedhypothesis}
For example, the pair estimate can be obtained by an eight-dimensional
Gaussian approximation with error $\delta_*$, followed by a comparison
with the independent marginal Gaussians with total variation error
$\chi$. We include in the exceptional count the ordinary and conjugate
angular diagonals, the radial duplicates, and $i=j$, unless they satisfy
the above pair estimate.

\begin{theorem}[Annular small-derivative count]\label{thm:T3}
Assume the \hyperref[hyp:mesh]{mesh and local multiplicity} bounds and the
\hyperref[hyp:jets]{jet approximation}, and denote by $B_N$ the number of
roots in $\mathcal A$ satisfying $|f_N'(\alpha)|<N^{3/2}/L$.
In the jet hypothesis use
\[
 u=\frac{N^{3/2}\tau/L+M_2\tau^2/2}{\sqrt N},\qquad
 v_0=\frac{R(N^{3/2}/L+M_2\tau)}{N^{3/2}}.
\]
We set
\[
 \mu_* =m\left(\frac{u^2v_0^2}{4\lambda^2}+\delta\right),
 \qquad V_*=b_*+\delta_*+\chi+2\delta.
\]
For every $t>0$,
\begin{equation}
 \Prob\{B_N>b_0+q_0(\mu_*+t)\}
       \leq p_0+p_2+\frac{m^2V_*}{t^2}. \label{eq:T3-general}
\end{equation}
\end{theorem}

\begin{proof}
The idea is to detect each counted root outside the excluded sectors at
its assigned mesh point, where both components of $W$ are small, and to
bound the number of mesh points with small $W$ by its mean and its
variance, which the jet approximation controls. We then bound $B_N$ by
Chebyshev's inequality and the cell bound.

Let $\alpha$ be a root in $B_N$ outside the excluded sectors, and let
$z_i$ be its assigned mesh point. On the derivative event
$\{\sup_D|f_N''|\leq M_2\}$, since $f_N(\alpha)=0$,
$|f_N'(\alpha)|<N^{3/2}/L$, $|z_i-\alpha|\leq\tau$, and the segment
$[\alpha,z_i]$ lies in $D$ by the
\hyperref[hyp:mesh]{mesh and local multiplicity} bounds, it follows from
Taylor's theorem with integral remainder along the segment that
\begin{align*}
 |f_N(z_i)|
 &=\left|f_N'(\alpha)(z_i-\alpha)
   +\int_{[\alpha,z_i]}(z_i-\zeta)f_N''(\zeta)\dd\zeta\right|\\
 &\leq |f_N'(\alpha)|\,|z_i-\alpha|+\frac{M_2}2|z_i-\alpha|^2
 \leq\frac{N^{3/2}\tau}{L}+\frac{M_2\tau^2}2,\\
 |f_N'(z_i)|
 &=\left|f_N'(\alpha)+\int_{[\alpha,z_i]}f_N''(\zeta)\dd\zeta\right|
 \leq|f_N'(\alpha)|+M_2|z_i-\alpha|
 \leq\frac{N^{3/2}}L+M_2\tau.
\end{align*}
Since $|z_i|\leq R$, dividing by the normalizations in $W(z_i)$ we
obtain
\[
 \frac{|f_N(z_i)|}{\sqrt N}
 \leq\frac{N^{3/2}\tau/L+M_2\tau^2/2}{\sqrt N}
 =u,\qquad
 \frac{|z_if_N'(z_i)|}{N^{3/2}}
 \leq\frac{R(N^{3/2}/L+M_2\tau)}{N^{3/2}}
 =v_0,
\]
so that the event $E_i$ occurs at every assigned point.

We set $Z=\sum_{i=1}^m\boldsymbol1_{E_i}$.
Let $\Sigma_i$ be the covariance matrix of $H_i$, whose eigenvalues are at
least $\lambda$, so that it has determinant at least $\lambda^4$. Hence
the density of $H_i$ is bounded by $(2\pi)^{-2}\lambda^{-2}$, since at
every $y\in\mathbb R^4$ it satisfies
\[
 (2\pi)^{-2}(\det\Sigma_i)^{-1/2}
   \exp\bigl(-\tfrac12\langle\Sigma_i^{-1}y,y\rangle\bigr)
 \leq(2\pi)^{-2}(\det\Sigma_i)^{-1/2}
 \leq(2\pi)^{-2}\lambda^{-2}.
\]
Since the product of the two complex disks has four-dimensional volume
$\pi^2u^2v_0^2$, by the \hyperref[hyp:jets]{jet approximation} we have
\begin{align*}
 g_i
 &\leq(2\pi)^{-2}\lambda^{-2}\,\pi^2u^2v_0^2
 =\frac{u^2v_0^2}{4\lambda^2},\\
 \E Z
 &=\sum_{i=1}^m\Prob(E_i)
 \leq\sum_{i=1}^m(g_i+\delta)
 \leq m\left(\frac{u^2v_0^2}{4\lambda^2}+\delta\right)
 =\mu_*.
\end{align*}

For an admissible ordered pair $(i,j)$, that is, one that obeys the pair
estimate of the jet approximation, since $\Prob(E_i),g_j\leq1$ we have
\begin{align*}
 |\Prob(E_i)\Prob(E_j)-g_ig_j|
 &=\bigl|\Prob(E_i)\bigl(\Prob(E_j)-g_j\bigr)
       +g_j\bigl(\Prob(E_i)-g_i\bigr)\bigr|\\
 &\leq\Prob(E_i)|\Prob(E_j)-g_j|
       +g_j|\Prob(E_i)-g_i|
 \leq2\delta,\\
 |\operatorname{Cov}(\boldsymbol1_{E_i},\boldsymbol1_{E_j})|
 &=|\Prob(E_i\cap E_j)-\Prob(E_i)\Prob(E_j)|\\
 &\leq|\Prob(E_i\cap E_j)-g_ig_j|+|g_ig_j-\Prob(E_i)\Prob(E_j)|
 \leq\delta_*+\chi+2\delta.
\end{align*}
On each of the at most $b_*m^2$ exceptional pairs we use the bound one,
and summing over all $m^2$ ordered pairs, including the diagonal, we
obtain
\begin{align*}
 \operatorname{Var}Z
 &=\sum_{i,j=1}^m
       \operatorname{Cov}(\boldsymbol1_{E_i},\boldsymbol1_{E_j})
 \leq\sum_{i,j=1}^m
       |\operatorname{Cov}(\boldsymbol1_{E_i},\boldsymbol1_{E_j})|\\
 &\leq b_*m^2+m^2(\delta_*+\chi+2\delta)
 =m^2V_*.
\end{align*}
Since $\E Z\leq\mu_*$, Chebyshev's inequality yields
\[
 \Prob(Z>\mu_*+t)
 \leq\Prob(Z-\E Z>t)
 \leq\frac{\operatorname{Var}Z}{t^2}
 \leq\frac{m^2V_*}{t^2}.
\]

On the two geometric good events, each cell receives at most $q_0$
roots, and the excluded sectors contain at most $b_0$ roots. On these
events and on the derivative event, by the first part of the proof every
counted root outside the excluded sectors is assigned to a retained point
$z_i$ at which $E_i$ occurs, and the cell bound of the
\hyperref[hyp:mesh]{mesh hypothesis}, applied to the cell of $z_i$, which
contains the roots assigned to $z_i$, shows that at most $q_0$ roots are
assigned to each such point. Hence $B_N\leq b_0+q_0Z$, so that
$B_N>b_0+q_0(\mu_*+t)$ implies $Z>\mu_*+t$ outside the geometric and
derivative exceptional events, and
\[
 \Prob\{B_N>b_0+q_0(\mu_*+t)\}
 \leq p_0+p_2+\Prob(Z>\mu_*+t)
 \leq p_0+p_2+\frac{m^2V_*}{t^2}.
\]
\end{proof}

\begin{corollary}[Polynomial parameters for the mesh bound]\label{cor:T3-powers}
Suppose the following bounds hold with fixed positive constants:
\begin{align*}
 &R\leq2,\qquad A\geq S,\qquad \tau=(ANL\sqrt{\log N})^{-1},\\
 &M_2=SN^{5/2}\sqrt{\log N},\qquad
 m\leq DNL^2\log N,\qquad q_0\leq C_0\log N.
\end{align*}
\[
 \delta\leq D_0N^{-\beta_1},\qquad
 V_*\leq VN^{-\beta},\qquad L=N^\ell.
\]
Then
\begin{equation}
 \mu_*\leq\frac{9D}{A^2\lambda^2}NL^{-4}
           +DD_0N^{1-\beta_1}L^2\log N. \label{eq:T3-mean}
\end{equation}
For a chosen $d>0$, impose the explicit deterministic thresholds
\[
 b_0\leq\tfrac12N^{1-d},\qquad
 C_0\log N\left[\frac{9D}{A^2\lambda^2}NL^{-4}
           +DD_0N^{1-\beta_1}L^2\log N\right]
       \leq\tfrac14N^{1-d}.
\]
Then
\begin{equation}
 \Prob(B_N>N^{1-d})
 \leq p_0+p_2+16C_0^2D^2V
         N^{-\beta+4\ell+2d}(\log N)^4. \label{eq:T3-power}
\end{equation}
The mean thresholds are eventual if
$d<4\ell$, $d<\beta_1-2\ell$, and $b_0=o(N^{1-d})$.
The last failure term is summable on $N=j^s$ whenever
\[
 s(\beta-4\ell-2d)>1.
\]
In particular $\beta=\beta_1=1/2$, $\ell=1/64$, $d=1/32$
give $B_N\leq N^{31/32}$.  The probability of the exceptional event, in
addition to $p_0+p_2$, is at most
\[
 16C_0^2D^2VN^{-3/8}(\log N)^4.
\]
\end{corollary}

\begin{proof}
We first substitute the parameters into the definitions of $u$, $v_0$
and $\mu_*$ in Theorem~\ref{thm:T3}, and then apply this theorem with a
deviation $t$ such that $b_0+q_0(\mu_*+t)\leq N^{1-d}$.

With the chosen parameters, and since $S\leq A$ and $R\leq2$, we have
\begin{align*}
 \frac{N^{3/2}\tau}L
 &=\frac{N^{1/2}}{AL^2\sqrt{\log N}},\qquad
 \frac{M_2\tau^2}2
 =\frac{SN^{1/2}}{2A^2L^2\sqrt{\log N}},\\
 u
 &=\left(\frac1A+\frac{S}{2A^2}\right)
                 L^{-2}(\log N)^{-1/2}
 \leq\frac3{2A}L^{-2}(\log N)^{-1/2},\\
 M_2\tau
 &=\frac SA\cdot\frac{N^{3/2}}L,\qquad
 v_0
 =\frac RL\left(1+\frac SA\right)
 \leq\frac4L,\\
 \frac{u^2v_0^2}{4\lambda^2}
 &\leq\frac1{4\lambda^2}
       \frac9{4A^2L^4\log N}\frac{16}{L^2}
 =\frac9{A^2\lambda^2}\frac{L^{-6}}{\log N}.
\end{align*}
Using also $m\leq DNL^2\log N$ and $\delta\leq D_0N^{-\beta_1}$, we have
\[
 \mu_*
 \leq DNL^2\log N
      \left(\frac9{A^2\lambda^2}\frac{L^{-6}}{\log N}
                   +D_0N^{-\beta_1}\right)
 =\frac{9D}{A^2\lambda^2}NL^{-4}
      +DD_0N^{1-\beta_1}L^2\log N.
\]
This proves \eqref{eq:T3-mean}.

We now choose the deviation $t$ in Theorem~\ref{thm:T3}. The
event $\{B_N>N^{1-d}\}$ is contained in the event of
\eqref{eq:T3-general} when $b_0+q_0(\mu_*+t)\leq N^{1-d}$. Since the
probability bound $m^2V_*/t^2$ decreases in $t$, we take the largest $t$
with $C_0t\log N\leq\frac14N^{1-d}$,
\[
 t=\frac{N^{1-d}}{4C_0\log N}.
\]
By the two deterministic thresholds, the mean bound \eqref{eq:T3-mean}
with $q_0\leq C_0\log N$, and this choice of $t$, we have
\begin{align*}
 b_0+q_0(\mu_*+t)
 &=b_0+q_0\mu_*+q_0t\\
 &\leq\frac12N^{1-d}
   +C_0\log N\left[\frac{9D}{A^2\lambda^2}NL^{-4}
           +DD_0N^{1-\beta_1}L^2\log N\right]\\
 &\quad+C_0t\log N
 \leq\frac12N^{1-d}+\frac14N^{1-d}+\frac14N^{1-d}
 =N^{1-d}.
\end{align*}
Hence, by \eqref{eq:T3-general}, $m\leq DNL^2\log N$,
$V_*\leq VN^{-\beta}$ and $L=N^\ell$, we have
\begin{align*}
 \Prob(B_N>N^{1-d})
 &\leq\Prob\{B_N>b_0+q_0(\mu_*+t)\}
 \leq p_0+p_2+\frac{m^2V_*}{t^2},\\
 \frac{m^2V_*}{t^2}
 &\leq D^2N^2L^4(\log N)^2\,VN^{-\beta}
           \frac{16C_0^2(\log N)^2}{N^{2-2d}}\\
 &=16C_0^2D^2VN^{-\beta}L^4N^{2d}(\log N)^4
 =16C_0^2D^2V N^{-\beta+4\ell+2d}(\log N)^4.
\end{align*}
We have thus proved \eqref{eq:T3-power}.
With $L=N^\ell$, the left-hand side of the second threshold equals
\begin{align*}
 &\frac{9C_0D}{A^2\lambda^2}N^{1-4\ell}\log N
 +C_0DD_0N^{1-\beta_1+2\ell}(\log N)^2\\
 &\qquad=\frac14N^{1-d}\left(\frac{36C_0D}{A^2\lambda^2}N^{d-4\ell}\log N
 +4C_0DD_0N^{d-\beta_1+2\ell}(\log N)^2\right).
\end{align*}
When $d<4\ell$ and $d<\beta_1-2\ell$, the factor in parentheses tends to
zero, and so the second threshold holds for all large $N$. The first
threshold holds for all large $N$ when $b_0=o(N^{1-d})$. With $N=j^s$,
the last failure term equals
$16C_0^2D^2Vj^{-s(\beta-4\ell-2d)}(s\log j)^4$, which is summable when
$s(\beta-4\ell-2d)>1$.
For the particular powers of the corollary, $4\ell=2d=1/16$, and the
three exponents, the conditions $d<4\ell$ and $d<\beta_1-2\ell$, and the
summability at $N=j^8$ read
\begin{gather*}
 1-4\ell=\frac{15}{16},\qquad
 1-\beta_1+2\ell=\frac{17}{32},\qquad
 -\beta+4\ell+2d=-\frac38,\\
 4\ell=\frac1{16}>\frac1{32}=d,\qquad
 \beta_1-2\ell=\frac{15}{32}>d,\qquad
 8(\beta-4\ell-2d)=3>1.
\end{gather*}
Since $1-d=31/32$, once the two thresholds hold, we obtain
$B_N\leq N^{31/32}$ outside an event of probability at most
$p_0+p_2+16C_0^2D^2VN^{-3/8}(\log N)^4$.
\end{proof}

\subsection[From a sparse subsequence of degrees to the full sequence]{From a sparse subsequence of degrees\\to the full sequence}

\begin{theorem}[Sparse-to-dense convergence]\label{thm:T5}
Let $s\geq1$ be fixed, $n_j=\lfloor j^s\rfloor$, and let $(X_n)$
be real random variables on a common probability space.  Suppose
$X_{n_j}/n_j\to c$ almost surely.  If there are deterministic
$\varepsilon_j\to0$ and events
$E_j$ such that
\[
 \sum_j\Prob(E_j)<\infty,\qquad
 \max_{n_j\leq n<n_{j+1}}|X_n-X_{n_j}|
       \leq\varepsilon_j n_j\quad\hbox{on }E_j^c,
\]
then $X_n/n\to c$ almost surely.
\end{theorem}
The numbers $\varepsilon_j$ need not be monotone.
In particular, the case $c=1/2$ yields the strong law needed for the
unweighted Rademacher model.

\begin{proof}
We compare each degree $n$ of the block $n_j\leq n<n_{j+1}$ with $n_j$,
and use the Borel--Cantelli lemma to discard the events $E_j$.
For every $J$, by the union bound we have
\[
 \Prob\Big(\bigcup_{j\geq J}E_j\Big)\leq\sum_{j\geq J}\Prob(E_j)\longrightarrow0
\]
as $J\to\infty$, since the series converges.
It follows that almost surely only finitely many $E_j$ occur, which is the
Borel--Cantelli lemma.
For $n_j\leq n<n_{j+1}$ we write
\[
 \frac{X_n}{n}-c
 =\frac{n_j}{n}\left(\frac{X_{n_j}}{n_j}-c\right)
   +\frac{X_n-X_{n_j}}n+c\left(\frac{n_j}n-1\right).
\]
Outside $E_j$, by the block bound $|X_n-X_{n_j}|\leq\varepsilon_jn_j$, it
follows that
\begin{align*}
 \left|\frac{X_n}{n}-c\right|
 &\leq\frac{n_j}{n}\left|\frac{X_{n_j}}{n_j}-c\right|
      +\frac{|X_n-X_{n_j}|}n+|c|\left|\frac{n_j}n-1\right|\\
 &\leq\frac{n_j}{n}\left|\frac{X_{n_j}}{n_j}-c\right|
      +\frac{\varepsilon_jn_j}n+|c|\frac{n-n_j}n\\
 &\leq\frac{n_j}{n}\left|\frac{X_{n_j}}{n_j}-c\right|
      +\varepsilon_j+|c|\frac{n_{j+1}-n_j}{n_j},
\end{align*}
where the last step uses $\varepsilon_j\geq0$, which holds on $E_j^c$
since there $0\leq|X_n-X_{n_j}|\leq\varepsilon_jn_j$.
Since $j^s-1<n_j\leq j^s$ for every $j$, we have, for $j\geq2$,
\[
 \frac{n_{j+1}}{n_j}
 <\frac{(j+1)^s}{j^s-1}
 =\frac{(1+1/j)^s}{1-j^{-s}}
 \longrightarrow1.
\]
Since also $n_{j+1}\geq n_j$, it follows that $n_{j+1}/n_j\to1$, and
hence the last ratio, which equals $n_{j+1}/n_j-1$, tends to zero.
As $n\to\infty$, the index $j$ with $n_j\leq n<n_{j+1}$ tends to
infinity. Almost surely, $X_{n_j}/n_j\to c$ and $E_j$ occurs for only
finitely many $j$, and then the three terms of the last bound tend
to zero, the first since $n_j/n\leq1$ and the second since
$\varepsilon_j\to0$.
This proves the assertion.
\end{proof}

\subsection{Deterministic choice of a vanishing block envelope}

\begin{lemma}[Deterministic diagonalization]\label{lem:deterministic-diagonalization}
Let $s\geq1$, $n_j=\lfloor j^s\rfloor$, and let $(X_n)$ be real
random variables on a common probability space.
Suppose that for each positive integer $k$,
\[
 \Prob\left\{\max_{n_j\leq n<n_{j+1}}
 \frac{|X_n-X_{n_j}|}{n_j}>a_k+b_k(j)\right\}\leq p_k(j),
\]
where $a_k\to0$, $b_k(j)\to0$ for fixed $k$, and
$\sum_jp_k(j)<\infty$ for fixed $k$.  Then there are deterministic
numbers $\varepsilon_j\to0$ and events $E_j$ with
$\sum_j\Prob(E_j)<\infty$ such that
\[
 \max_{n_j\leq n<n_{j+1}}\frac{|X_n-X_{n_j}|}{n_j}\leq\varepsilon_j
 \quad\hbox{on }E_j^c.
\]

In particular, if $X_{n_j}/n_j\to c$ almost surely, then
$X_n/n\to c$ almost surely.
\end{lemma}

\begin{proof}
The idea is to let the index $k$ grow slowly with $j$, along a
deterministic diagonal $k(j)\to\infty$ chosen so that the probabilities
of the exceptional events along the diagonal are still summable, and then
to apply Theorem~\ref{thm:T5}.
Choose increasing deterministic integers
$J_k$, with $J_k\geq k$ and $J_k>J_{k-1}$, such that for all
$j\geq J_k$ we have $b_k(j)\leq1/k$ and
$\sum_{j\geq J_k}p_k(j)\leq2^{-k}$, which is possible since $b_k(j)\to0$
and $\sum_jp_k(j)<\infty$.
Choosing the least admissible integer makes each $J_k$ deterministic,
and this definition does not require a numerical bound for $J_k$.
For $j\geq J_1$, we set
\[
 k(j)=\max\{k:J_k\leq j\},\qquad
 \varepsilon_j=a_{k(j)}+\frac1{k(j)}.
\]
Since $J_k\geq k$ and the $J_k$ increase strictly, $k(j)$ is finite
and tends to infinity, and therefore $\varepsilon_j\to0$. Let $E_j$ be
the event that the normalized block increment exceeds $\varepsilon_j$,
and for $j<J_1$ take $E_j=\Omega$ and choose $\varepsilon_j$ arbitrarily.
For $j\geq J_1$ we have $J_{k(j)}\leq j$, hence $b_{k(j)}(j)\leq1/k(j)$
and $\varepsilon_j\geq a_{k(j)}+b_{k(j)}(j)$. By the hypothesis with
$k=k(j)$ we therefore have $\Prob(E_j)\leq p_{k(j)}(j)$, and these bounds
satisfy
\[
 \sum_{j\geq J_1}p_{k(j)}(j)
 \leq\sum_{k\geq1}\sum_{j\geq J_k}p_k(j)
 \leq\sum_{k\geq1}2^{-k}
 =1,
\]
where the first step holds since each $j\geq J_1$ contributes the single
nonnegative term $p_{k(j)}(j)$, with $j\geq J_{k(j)}$, which appears in the
inner sum for $k=k(j)$.
Since only finitely many $j<J_1$ remain, it follows that
$\sum_j\Prob(E_j)<\infty$.
Finitely many initial degrees do not affect the limit, and the final
assertion follows from Theorem~\ref{thm:T5}.
The conclusion holds for arbitrary dependence between degrees,
blocks, and exceptional events.
\end{proof}

\subsection{The mean for coefficients in \texorpdfstring{$\{0,1\}$}{\{0,1\}}}
\begin{remark}[The mean for coefficients in $\{0,1\}$]\label{rem:mean-newman}
Let $K\ge2$, $N\ge2K$, $|r-1|\le K/N$ and
$\operatorname{dist}(\theta,2\pi\mathbb Z)>0$, let
$\sigma_N(r)^2=\sum_{k=0}^Nr^{2k}$, as in Section~\ref{sec:radial}, and let
$D_N(z)=\sum_{k=0}^Nz^k$. Then
\[
 \frac{|D_N(re^{i\theta})|}{\sigma_N(r)}
 \le\frac{10e^{4K}}{\sqrt N\operatorname{dist}(\theta,2\pi\mathbb Z)}.
\]
Indeed, $D_N(re^{i\theta})=\{1-(re^{i\theta})^{N+1}\}/(1-re^{i\theta})$.
For the numerator, $r\le1+K/N$ gives
$|1-(re^{i\theta})^{N+1}|\le1+(1+K/N)^{N+1}\le1+e^{2K}\le2e^{2K}$.
For the denominator,
$|1-re^{i\theta}|^2=(1-r)^2+4r\sin^2(\theta/2)\ge4r\sin^2(\theta/2)$,
and $r\ge1-K/N\ge1/2$ together with
$|\sin(\theta/2)|\ge\operatorname{dist}(\theta,2\pi\mathbb Z)/\pi$, by the
concavity of the sine on $[0,\pi/2]$, gives
$|1-re^{i\theta}|\ge(\sqrt2/\pi)\operatorname{dist}(\theta,2\pi\mathbb Z)$.
For the normalization, every term of $\sigma_N(r)^2$ is at least $(1-K/N)^{2N}\ge e^{-4K}$, by $\log(1-u)\ge-2u$ for
$0\le u\le1/2$, so that $\sigma_N(r)\ge\sqrt Ne^{-2K}$.
Combining the three bounds and $\sqrt2\pi<10$ gives the claim.
\end{remark}
